\section{Cohomological Overdetermination of the Residue $1/8$}
\label{sec:surdetermination_un_huitieme}

The semi-classical PT density differs from the Riemann--von Mangoldt
density by a constant $1/8$ (\S\ref{sec:densite_semi_classique}).
\citeauthor{BerryKeating1999} \citeyear{BerryKeating1999} identify
this constant as the fine Maslov phase of the corners of the double
barrier, without dynamical interpretation. PT provides a \emph{triple
cohomological identification} that converges, overdetermined by the
coherence of the three readings.

\subsection{K2 --- Algebraic identity: $1/8 = \sper^{2}/2 = \lambdaH$}
\label{sec:K2}

The scalar Higgs self-coupling in the PT sector is derived in the
monograph (ch.~15, depth $d = 3$) \cite{PT_Foundations}:
\begin{equation}
\lambdaH \;=\; \dfrac{\sper^{2}}{2}.
\label{eq:lambdaH}
\end{equation}
At $\sper = 1/2$, $\lambdaH = 1/8$. This identity is structurally
tied to the two-branch bifurcation $\qplus/\qminus$: $\lambdaH$ is
the square of the bifurcation parameter divided by the number of
scalar branches.

\paragraph{Associated algebraic identifications.}
Two other algebraic identifications coincide:
\[
\dfrac{1}{8} \;=\; \bigg(\dfrac{1}{2}\bigg)^{3}
\quad \text{(cube of the scalar unit, \cite{PT_M5})},
\]
\[
\dfrac{1}{8} \;=\; \dfrac{1}{2 N_{\mathrm{gen}} + 2}
\quad \text{(at $N_{\mathrm{gen}} = 3$, exact)}.
\]

\paragraph{Rejected naive candidates.}
Several candidates are rejected by numerical incompatibility:
$1/N_{\mathrm{Weyl}} = 1/15 \ne 1/8$,
$\chi(\Sigma_{\mathrm{pers}})/Z = -26/Z \ne 1/8$,
$\gammathree \gammafive \gammaseven \approx 0.335 \ne 1/8$,
$\alpha_{\mathrm{bare}} \approx 1/136 \ne 1/8$. The $1/8$ is not a
topological invariant of the persistence surface, but a scalar
invariant of the bifurcation $\qplus/\qminus$.

\subsection{K3 --- Cohomological identity: $1/8 = c_1/N_{\mathrm{corners}}$}
\label{sec:K3}

The relevant cohomological structure is constructed in the PT corpus.
The doubled Fisher manifold
\begin{equation}
\mathcal{M}_m^{\mathrm{db}} \;=\; \mathcal{M}_m \times \mathcal{M}_m, \qquad
g(p,q) \;=\; g_{\mathrm{Fisher}}(p) \,\oplus\, g_{\mathrm{Fisher}}(p),
\label{eq:Mdoub}
\end{equation}
(the \emph{same} base point $p$ in both blocks --- canonical
Dombrowski 1962 construction on the tangent bundle of a Hessian
manifold), equipped with the complex structure $J_{\mathrm{PT}}(u, v) = (-v, u)$
is a K\"ahler manifold (theorem~\texttt{kahler\_fisher},
\cite{PT_M5}). Its K\"ahler potential is the negative Shannon
entropy $K(p) = \sum_k p_k \log p_k$ (canonical Hessian potential of
the Fisher metric, Amari--Nagaoka 2000). The chiral projectors
\begin{equation}
P_{\pm} \;=\; \tfrac{1}{2}(I \mp i J_{\mathrm{PT}})
\label{eq:Pchir}
\end{equation}
decompose $H^{*}(BG_m; \Cplx)$ into $\pm i$ eigenspaces and are
identified with the spinor projector $\qplus/\qminus$ of the sieve
(\cite{PT_M5}, \texttt{def:chiral\_projectors}).

\begin{theoreme}[First Chern class of the PT spinor bundle, \texttt{berry\_3pi}]
\label{thm:berry_3pi}
The first Chern class of the spinor bundle over
$\mathcal{M}_m^{\mathrm{db}}$ equals
\begin{equation}
c_1 \;=\; \sper \;=\; \dfrac{1}{2}.
\label{eq:c1}
\end{equation}
% ERRATUM E-PM-8 (2026-08-15): qualification of the SECOND statement. It was
% given with no conditionality marker inside a "theorem" environment, whereas
% the berry_3pi theorem of the monograph (app. P, sec:appP_C6_promotion) carries
% [COND]. No mathematical content is modified.
The holonomy of the $N_{\mathrm{gen}}$-fold cover is
$\exp(2\pi i\,c_1\,N_{\mathrm{gen}}) = \exp(3\pi i)$ (origin of the
\texttt{berry\_3pi} theorem of the monograph, app.~P, \S~C7) --- this
second statement is \textbf{[COND]}: it assumes $N_{\mathrm{gen}} = 3$
as an explicit model premise, which the colour-counting theorem does not
establish, and the monograph's \texttt{berry\_3pi} theorem is itself
\textbf{[COND]}. Only $c_1 = 1/2$ above is unconditional.
\end{theoreme}

The algebraic sub-statement $c_1 = 1/2$ is \textbf{[Thm]}
unconditional; full APS rigour of the calculation
$1/8 = c_1/N_{\mathrm{corners}}$ on the cornered manifold remains
partially open (\S\ref{sec:limites_ouvertures}), status
\textbf{[Der]/[partial]}.

\begin{proposition}[Cohomological computation of the residue $1/8$]
\label{prop:c1_sur_N}
In the canonical phase space $(u, p_u)$, the BK double barrier is a
symplectic quarter-plane
$Q_\star = \{u \ge u_\star,\, p_u \ge p_\star,\, u_\star\,p_\star = 2\pi\}$.
The boundary admits two edges meeting at the corner. With the Weyl
symmetrisation and the two mirror branches identified by the
antiperiodicity $T_3$ (\S\ref{sec:bord_antiperiodique}), one obtains
\begin{equation}
N_{\mathrm{corners}} \;=\; 4
\label{eq:Ncoins}
\end{equation}
distinct corners on the complete cell. The Berry phase contributed by
each corner, measured on the quarter-turn joining the two adjacent
edges, is
\begin{equation}
\gamma_{\mathrm{corner}} \;=\; -\dfrac{2\pi\,c_1}{N_{\mathrm{corners}}} \;=\; -\dfrac{\pi}{4},
\label{eq:gamma_corner}
\end{equation}
and the associated Bohr--Sommerfeld density is
\begin{equation}
\boxed{\;N_{\mathrm{corner}} \;=\; \dfrac{\gamma_{\mathrm{corner}}}{2\pi} \;=\; -\dfrac{c_1}{N_{\mathrm{corners}}} \;=\; -\dfrac{1}{8}.\;}
\label{eq:c1_sur_N}
\end{equation}
\end{proposition}

\begin{proof}[Sketch]
The symplectic quarter-plane admits 2 edges ($u = u_\star$ and
$p_u = p_\star$) that meet at the corner
$(u_\star, p_\star) = (\sqrt{2\pi}, \sqrt{2\pi})$. The Weyl
symmetrisation (involution $(u, p_u) \leftrightarrow (p_u, u)$)
creates the mirror sector, doubling the edges and producing two
additional corners. The antiperiodicity $\theta = \pi$
(proposition~\ref{prop:antiperiodicite}) identifies two pairs of
corners into two, giving $N_{\mathrm{corners}} = 4$ in total (without
the identification, one would have $N = 8$ and $c_1/N = 1/16$,
contradicting the observed residue~\eqref{eq:residu_un_huitieme}).
The Berry phase per corner is computed by parallel transport on the
inner quarter-turn: in local polar variables centred on the corner,
$\gamma_{\mathrm{corner}} = \oint A = -2\pi\,c_1/N_{\mathrm{corners}} = -\pi/4$
with $c_1 = 1/2$. The passage to the Bohr--Sommerfeld density divides
by $2\pi$: $N_{\mathrm{corner}} = -1/8$. Numerical verification:
machine precision $1.68 \cdot 10^{-14}$ via the script
\texttt{pt\_k3\_cohomological\_maslov.py}.
\end{proof}

The identification $1/8 = c_1/N_{\mathrm{corners}}$ therefore yields
the Berry--Keating $-\pi/8$ per corner, but \emph{derived} from the
PT cohomological structure, not as an opaque geometric artefact.

\subsection{Berry phase per corner: $1/8 = \sper/4$}
\label{sec:phase_berry}

At the equator of the Bloch sphere ($\theta = \pi/2$), the Berry
connection of a spin-$\sper$ is constant:
$A_\varphi = -\sin^{2}(\pi/4) = -1/2 = -\sper$ per unit azimuthal
angle. Over a complete equatorial cycle, the Berry phase is
$\gamma_1 = -\pi = -2\pi\,\sper$. Divided by the four
corners~\eqref{eq:Ncoins}:
\begin{equation}
\dfrac{\gamma_1}{N_{\mathrm{corners}}} \;=\; -\dfrac{\pi}{4} \;=\; -2\pi \cdot \dfrac{\sper}{4},
\label{eq:s_sur_4}
\end{equation}
that is, $-1/8$ per corner in area count. The corresponding algebraic
identity is
\begin{equation}
\dfrac{1}{8} \;=\; \dfrac{\sper}{4}.
\label{eq:un_huitieme_sur_4}
\end{equation}

\subsection{Overdetermination by coherence}
\label{sec:surdetermination_coherence}

The three identifications converge to the same number at $\sper = 1/2$:
\begin{equation}
\boxed{\;
\dfrac{1}{8} \;=\; \underbrace{\dfrac{\sper^{2}}{2}}_{\text{K2, Higgs}} \;=\; \underbrace{\dfrac{c_1}{N_{\mathrm{corners}}}}_{\text{K3, Chern}} \;=\; \underbrace{\dfrac{\sper}{4}}_{\text{Berry per corner}}.\;}
\label{eq:triple_identification}
\end{equation}
The algebraic coherence between K2 and the Berry phase per corner
reduces to
\begin{equation}
\dfrac{\sper^{2}}{2} \;=\; \dfrac{\sper}{4} \quad\Longleftrightarrow\quad \sper\,(2\sper - 1) \;=\; 0.
\label{eq:coherence_K2_Berry}
\end{equation}
The only solutions are $\sper = 0$ (trivial, no active parameter,
excluded by $\Tone$) and $\sper = 1/2$ (theorem $\Tone$). \textbf{The
coherence of the three cohomological readings overdetermines
$\sper = 1/2$ modulo the assumed $\Tone$ axiom}; it does not derive
$\sper = 1/2$ independently of $\Tone$. It is a structural signature
of $\Tone$ from the HP-PT zone of the corpus, not a substitute for
$\Tone$.

\subsection{Index-theorem reading}
\label{sec:indice_APS}

The final assembly is interpreted as a formal
Atiyah--Patodi--Singer index theorem \cite{APS1975} on the cornered
manifold $\square$:
\begin{equation}
\mathrm{ind}\!\big(D_+|_\square\big) \;=\; \int_{\square} \mathrm{ch}(E)\,\mathrm{Td}(\square) \,+\, \dfrac{1}{2}\,\big(\eta_{\partial} + \dim \ker D_\partial\big) \,+\, \sum_{\mathrm{coins}} \dfrac{\delta_{\mathrm{coin}}}{2}.
\label{eq:APS}
\end{equation}
With $c_1(E) = 1/2$ and $\delta_{\mathrm{coin}} = \pi/2$ for a right
corner, the per-corner contribution is $\pi/2 / (2 \cdot 2\pi) = 1/8$,
in agreement with the direct Berry calculation. The equivalent
identification via Chern--Simons:
$\exp(2\pi i\,c_1/N_{\mathrm{corners}}) = \exp(i\pi/4)$, whose argument
$\pi/4$ gives in density $1/8$ per corner.

Formal rigour in the strict APS sense (beyond the area calculation
and numerical verification) remains to be developed
(\S\ref{sec:limites_ouvertures}). The cohomological mechanism is
nonetheless explicitly identified.
